\section{Exercise 1}
\textit{Consider the two lines $l$ and $m$ in $\mathbb{R}^3$ given by:}
\[l = \{(1 +2t, -1+4t,2+3t) \in \mathbb{R}^3 | t \in \mathbb{R}\}\]
\[m = \{(1-s,1+2s,4s)\in \mathbb{R}^3 | s \in \mathbb{R}\}\]

\begin{itemize}
    \item Argue i and m are NOT parallel, using the cross-product
    \item Write down a function of (t,s) that outputs the square distance between a point on l and a point m and use this function to find the minimum distance between the two lines
\end{itemize}

l and m can be written down as a starting point and a direction vector:
\[l(t) = 
\begin{pmatrix}
    1\\
    -1\\
    2
\end{pmatrix}
+ t
\begin{pmatrix}
    2\\
    4\\
    3
\end{pmatrix},
\qquad
m(s) = 
\begin{pmatrix}
    1\\
    1\\
    0
\end{pmatrix}
+ s 
\begin{pmatrix}
    -1\\
    2\\
    4
\end{pmatrix}
\]
By taking the cross product of the direction vectors, if this is $\neq 0$ then they are nonparallel

\[
l'(t) = \begin{pmatrix}
    2\\ 4\\ 3
\end{pmatrix}, \qquad
m'(s) = \begin{pmatrix}
    -1\\2\\4
\end{pmatrix}
\]
\[
l'(t) \times m'(s) = 
\begin{vmatrix}
    i&j&k\\
    2&4&3\\
    -1&2&4
\end{vmatrix}
\Rightarrow
\begin{pmatrix}
    4\cdot4-2\cdot3\\
    2\cdot4-(-1)\cdot3\\
    2\cdot2-(-1)\cdot4
\end{pmatrix}
=
\begin{pmatrix}
    10\\
    11\\
    8
\end{pmatrix}
\]
Since $\begin{pmatrix} 10\\11\\8 \end{pmatrix}$ is
not a zero vector, then they are not parallel.

To solve this, it is required to build a distance function with two inputs.
This distance function consists of a vector that points from one vector to the other, 
then finding the distance, then squaring it.

\[d(t,s) = l(t) - m(s) = 
\begin{pmatrix}
    (1+2t) - (1-s)\\
    (-1+4t) - (1+2s)\\
    (2+3t)-4s
\end{pmatrix}\]

The squared distance is then:
\[d(t,s) = |l(t) - m(s)|^2\]
\[\Downarrow\]
\[d(t,s) = (2t+s)^2 + (-2+4t-2s)^2 + (2+3t-4s)^2\]
To find a minimum, the goal is to find a minimum, where small changes in either t or s, is zero, such that:
\[
\frac{\partial d(t,s)}{\partial t} = 0, \qquad
\frac{\partial d(t,s)}{\partial s} = 0
\]
The distance expression, can then be partial differentiated to seek this minimum.
% This requires the chain rule: $(g(x))^2 = 2\left(g(x)\right)g'(x)$.
\[
\frac{\partial d(t,s)}{\partial t} 
\begin{pmatrix}
    (2t+s)^2 \\
    (-2+4t-2s)^2 \\
    (2+3t-4s)^2
\end{pmatrix}
= 
\begin{pmatrix}
    2(2t+s)(2)\\
    2\left((-2)+4t-2s\right)(4)\\
    2(2+3t-4s)(3)
\end{pmatrix}
\Rightarrow
\begin{pmatrix}
    8t+4s\\
    32t-16s-16\\
    18t-24s+12
\end{pmatrix}
\]
This expression is then summed together
\[
\frac{\partial d(t,s)}{\partial t} =
8t+4s + 32t-16s-16 + 18t-24s+12 = 58t-36s-4
\]
The same is done for the partial derivative of d with respect to s.
\[
\frac{\partial d(t,s)}{\partial s} 
\begin{pmatrix}
    (2t+s)^2 \\
    (-2+4t-2s)^2 \\
    (2+3t-4s)^2
\end{pmatrix}
= 
\begin{pmatrix}
    2(2t+s)(1)\\
    2\left((-2)+4t-2s\right)(-2)\\
    2(2+3t-4s)(-4)
\end{pmatrix}
\Rightarrow
\begin{pmatrix}
    4t+2s\\
    -16t+8s+8\\
    -24t+32s-16
\end{pmatrix}
\]
Summed together:
\[
\frac{\partial d(t,s)}{\partial s} =
-36t+42s-8
\]
Earlier it was stated the intention is to find the place where the
gradient of the partials is equal to zero, thus, these two equations can
be set equal to zero and placed as a system of two equations with
two unknown, to solve for t and s.
\[
\begin{cases}
    58t-36s=4 \\
    -36t+42s=8
\end{cases}
\]
Using a CAS tool, this can be used to isolate s in one equation
and substitute that into the other yielding:
\[
\begin{cases}
    t = \frac{2}{5}\\
    s = \frac{8}{15}
\end{cases}
\]
These value are substituted into the expression for the squared
distance function, and then taken the square root of that expression
to yield the minimum distance, using a CAS tool:
\[\sqrt{d(t,s)} = \sqrt{d\left(\frac{2}{5}, \frac{8}{15}\right)} \approx 2.251 \]